\documentclass[11pt]{article}

\usepackage[margin=1in]{geometry}
\usepackage{amsmath,amssymb,amsthm}
\usepackage{graphicx}
\usepackage{microtype}
\usepackage[colorlinks=true,linkcolor=blue,citecolor=blue,urlcolor=blue]{hyperref}

\emergencystretch=2em

\newtheorem{theorem}{Theorem}
\newcommand{\C}{\mathbb{C}}
\newcommand{\N}{\mathbb{N}}
\newcommand{\Hil}{\mathcal{H}}
\newcommand{\Alg}{\mathcal{A}}
\newcommand{\Sd}{\mathrm{Sd}}
\newcommand{\Tr}{\operatorname{Tr}}
\newcommand{\spec}{\operatorname{spec}}

\title{A Diagonal Counterexample to ``Dimension Spectrum Encoded by \(D\)''}
\author{Run \texttt{fa\_banach\_001}, lane 16}
\date{2026-07-03}

\begin{document}
\maketitle

\section*{Status}

\textbf{Claim status: candidate counterexample, likely valid.}
This packet addresses Chapter 5, Problem 5 of Eckstein--Iochum
\cite{EcksteinIochum2019}.  The problem asks whether the dimension spectrum
of a spectral triple is always contained in the union of the pole sets of the
\(D\)-only zeta functions \(\zeta_{|D|^{-k},D}\).  The construction below
shows that this is false for regular spectral triples under the broad
definitions used in the source.

\section*{Source Excerpt}

The relevant question appears in Chapter 5, page 123 of the source PDF.

\begin{center}
\includegraphics[width=0.92\linewidth]{figures/open_problem_crop.png}
\end{center}

\section*{Transcribed Open Question}

The source observes that, in all worked-out examples known there,
\[
  \Sd \subset \bigcup_{k\in\N} \mathcal P(\zeta_{|D|^{-k},D}),
\]
and then asks whether the whole dimension spectrum is actually encoded in the
operator \(D\), or whether this is only a property of the worked-out examples.

\section*{Proof Intuition}

The spectrum of \(D\) controls the scalar zeta functions
\(\Tr |D|^{-s-k}\), but the algebra can insert bounded diagonal coefficients
into the trace.  If those coefficients are \(n^i\), then the corresponding
zeta function is \(\sum n^{i-s}=\zeta(s-i)\), whose pole is shifted from
\(1\) to \(1+i\).  Since the inserted diagonal unitary commutes with \(D\),
regularity is automatic and the operator is order zero.  Thus the algebra,
not just \(D\), contributes a genuinely new pole.

\begin{theorem}
There is a regular 1-dimensional spectral triple with simple dimension
spectrum for which the minimal dimension spectrum is not contained in
\[
  \bigcup_{k\in\N_0}\mathcal P(\zeta_{|D|^{-k},D}).
\]
\end{theorem}

\begin{proof}
Let \(\Hil=\ell^2(\N)\) with standard basis \((e_n)_{n\geq 1}\), and define
\[
  D e_n = n e_n .
\]
Then \(D\) is positive self-adjoint with compact resolvent.  Let \(U\) be the
diagonal unitary
\[
  U e_n = n^i e_n = e^{i\log n}e_n ,
\]
and let \(\Alg\) be the concrete unital \(*\)-subalgebra of
\(\mathcal B(\Hil)\) generated by \(U\).  Thus the representation of
\(\Alg\) on \(\Hil\) is faithful by definition.

For every \(a\in\Alg\), \(a\) is a bounded diagonal operator and
\([D,a]=0\).  Hence \((\Alg,\Hil,D)\) is a spectral triple, and it is regular
because all iterated commutators with \(|D|\) vanish on \(\Alg\) and on
\([D,\Alg]\).  Its metric dimension is \(1\), since
\[
  \Tr |D|^{-s}=\sum_{n=1}^\infty n^{-s}=\zeta(s)
\]
converges exactly for \(\operatorname{Re}s>1\).

The source's abstract pdo algebra is generated by \(\Alg\), \(D\), and
\(|D|\), up to smoothing remainders.  In the present example these generators
commute.  Therefore every finite principal pdo term is a finite linear
combination of operators \(U^m |D|^j |D|^{-p}\), with \(m\in\mathbb Z\) and
\(j,p\in\N_0\).  For an order-zero term, \(p-j\geq0\), and
\[
  \Tr\!\left(U^m |D|^{j-p}|D|^{-s}\right)
  =
  \sum_{n=1}^\infty n^{im} n^{-(s+p-j)}
  =
  \zeta(s+p-j-im).
\]
Thus all such principal terms have meromorphic continuation with at most
simple poles in the discrete set
\[
  \{\,1-k+im : k\in\N_0,\ m\in\mathbb Z\,\}.
\]
Smoothing remainders contribute holomorphic functions after choosing the
remainder order arbitrarily negative, so the triple has a simple dimension
spectrum, and the minimal dimension spectrum contains the poles that actually
occur.

Now take the order-zero pdo \(T=U\in\Alg\).  Its zeta function is
\[
  \zeta_{U,D}(s)
  =
  \Tr(U|D|^{-s})
  =
  \sum_{n=1}^\infty n^{i-s}
  =
  \zeta(s-i),
\]
which has a simple pole at \(s=1+i\).  Hence \(1+i\) belongs to the minimal
dimension spectrum.

On the other hand, the \(D\)-only probe functions are
\[
  \zeta_{|D|^{-k},D}(s)
  =
  \Tr |D|^{-k}|D|^{-s}
  =
  \zeta(s+k),
\]
whose only pole is \(s=1-k\), real, for \(k\in\N_0\).  Therefore
\[
  1+i \notin \bigcup_{k\in\N_0}\mathcal P(\zeta_{|D|^{-k},D}),
\]
although \(1+i\) is in the dimension spectrum.  This disproves the proposed
general inclusion.
\end{proof}

\section*{Verification Notes}

The construction satisfies the source's Definition 1.1 of spectral triple,
Definition 1.6 of regularity, and Definition 1.22 of dimension spectrum.  The
only point to flag for review is scope: because \([D,\Alg]=0\), this is a
degenerate noncommutative metric geometry.  That degeneracy is not excluded in
the problem as stated, and the example is intentionally designed to test the
unrestricted formulation.

The script \texttt{code/check\_poles.py} lists the visible pole locations for
the \(D\)-only probes and for the algebra element \(U\).

\section*{Novelty and Search Bounds}

Run-index search was performed for \texttt{1902.05306}, the title, and close
phrases around \texttt{spectral action}, \texttt{noncommutative geometry},
\texttt{spectral triple}, and \texttt{dimension spectrum}; no direct duplicate
was found.  Web searches for close variants of "dimension spectrum encoded in
D", "whole dimension spectrum encoded in the operator D", and "diagonal
spectral triple dimension spectrum Dirichlet series" did not reveal a later
explicit answer.

\section*{Limitations}

This packet refutes the unrestricted general-fact formulation.  It does not
settle a strengthened question for nondegenerate, irreducible, geometric, or
reconstruction-compatible spectral triples.  In such narrower classes, the
phenomenon observed by the source may still be meaningful.

\section*{Human Review Recommendation}

Check that the source's \(\Psi^0(\Alg)\) convention indeed includes the
algebra element \(U\) as an order-zero pdo.  Once that is accepted, the shifted
zeta pole \(\zeta(s-i)\) gives the counterexample immediately.

\begin{thebibliography}{9}
\bibitem{EcksteinIochum2019}
Michal Eckstein and Bruno Iochum,
\emph{Spectral Action in Noncommutative Geometry},
arXiv:1902.05306, 2019.

\bibitem{ConnesMoscovici}
Alain Connes and Henri Moscovici,
\emph{The local index formula in noncommutative geometry},
Geom. Funct. Anal. 5 (1995), 174--243.
\end{thebibliography}

\end{document}
